% Deutsche Parallelfassung zu foundations/operators.tex

\chapter{Typisierte Operatoren}
\label{chap:operators}

Die überarbeitete Krümmung hat einen überarbeiteten intrinsischen Zustand
erzeugt. Die nächste Frage lautet nicht, welche Begriffe existieren, sondern
welche Handlung jede folgende Konsequenz erzeugt. AIM repräsentiert solche
Handlungen als typisierte Operatoren.

\begin{definition}[Operator]
\label{def:operator}
Ein Operator ist eine Abbildung
\[
\mathcal O_c:\mathcal A\rightarrow\mathcal B,
\]
deren Eingabetyp, Ausgabetyp, anwendbare Qualifikation \(c\), Annahmen und
Ergebnisprovenienz deklariert sind.
\end{definition}

Der Index wird gezeigt, wenn die Qualifikation sonst mehrdeutig wäre. Ein
Algorithmus kann einen Operator implementieren; der ingenieurtechnische
Vertrag wird aber nicht allein durch die Implementierung definiert.

\begin{principle}[Zustands--Operator-Beziehung]
\label{princ:state-operator-relation}
Operatoren bewerten, transformieren oder verknüpfen Engineering States, ohne
selbst Bestandteile des Engineering Objects zu werden.
\end{principle}

\begin{proposition}
\label{prop:operator-explicit-dependencies}
Ingenieurtechnische Abhängigkeiten, welche die Bedeutung, Domäne oder den
Status eines Zustands ändern, müssen ausdrücklich durch Operatoren oder
Beziehungen repräsentiert werden, statt in einer undifferenzierten
Zustandsbeschreibung verborgen zu sein.
\end{proposition}

\section{Den bearbeiteten Zustand rekonstruieren}

Der geometrische Rekonstruktionsoperator kann schematisch geschrieben werden
als
\[
\mathcal D_{I,s_0}:
(\mathrm{pose2}(s_0),\kappa)
\longrightarrow \mathrm{pose2}|_I.
\]

\paragraph{Ausführungsvertrag.}
Eingabe: eine Anfangspose, eine Krümmungsfunktion und ein Intervall. Ausgabe:
die abgeleitete Zustandstrajektorie. Erforderliche Qualifikationen:
longitudinale Domäne und Randbedingungen. Anwendbarkeit: Die
Rekonstruktionsaxiome und Regularitätsannahmen müssen gelten. Bewahrte
Identität: das Alignment, dem der bearbeitete Zustand zugeordnet ist.
Provenienz: Änderung, Anfangszustand und Rekonstruktionsverfahren. Der
Operator kann weder die Zulässigkeit bestimmen noch entscheiden, ob der
bearbeitete Zustand einen freigegebenen Zustand ersetzen soll.

\section{Den Zustand metrisch realisieren}

\begin{definition}[Operator der metrischen Realisierung]
\label{def:metric-realization-operator}
\[
\mathcal G_c:\mathcal X\rightarrow\mathcal X_{\mathrm{metric}}.
\]
\end{definition}

\paragraph{Ausführungsvertrag.}
Eingabe: der intrinsische Zustand. Ausgabe: seine metrische Interpretation.
Erforderliche Qualifikationen: metrisches Modell, Einbettungsdaten und
anwendbare Konventionen. Anwendbarkeit: Der Kontext muss die Zustandsdomäne
und die angeforderten Größen abdecken. Bewahrte Identität: Engineering Object
und Zustandsrolle. Provenienz: intrinsische Quelle plus Realization Context.
Der Operator kann weder physisch bauen, eine Beobachtung erzeugen noch eine
Entscheidung auswählen.

\begin{definition}[Frame-Operator]
\label{def:frame-operator}
Der Frame-Operator konstruiert das von der anwendbaren Realisierung benötigte
lokale Frame:
\[
\mathcal F_c:\mathcal X_{\mathrm{metric}}\rightarrow\mathcal F.
\]
\end{definition}

Seine Ausgabe sind abgeleitete Orientierungsdaten. Er bewahrt die Identität
und erbt Zustands- und Realisierungsprovenienz; er kann nicht bestimmen,
welche Frame-Konvention ein Projekt verwenden soll.

\section{Gleis- und World-Beschreibungen verknüpfen}

\begin{definition}[Track-to-World-Operator]
\label{def:track-to-world-operator}
\[
\mathcal P_{\mathrm{tw},c}:
\Omega_{\mathrm{track}}\rightarrow\Omega_{\mathrm{world}}.
\]
\end{definition}

Er empfängt Gleiskoordinaten und ein realisiertes Alignment-Frame und erzeugt
einen World-Space-Punkt. Seine Anwendbarkeit erfordert ein gültiges lokales
Frame über der angeforderten Domäne. Das Ergebnis bewahrt das referenzierte
Objekt und trägt die Realisierungsprovenienz; es kann keinen physischen oder
beobachteten Zustand begründen.

\begin{definition}[World-to-Track-Operator]
\label{def:world-to-track-operator}
\[
\mathcal P_{\mathrm{wt},c}:
\Omega_{\mathrm{world}}\rightarrow\Omega_{\mathrm{track}}.
\]
\end{definition}

Er empfängt einen World-Space-Punkt und ein anwendbares realisiertes Alignment
und erzeugt eine mögliche gleisbezogene Interpretation. Existenz und
Eindeutigkeit hängen von Suchdomäne und Geometrie ab. Seine Provenienz muss
World-Quelle, Realisierung und Auswahlregel bewahren. Er kann weder beweisen,
dass ein Vermessungspunkt das Alignment beobachtet, noch ohne anwendbare Regel
zwischen mehrdeutigen Projektionen wählen.

\begin{definition}[Koordinatentransformationsoperator]
\label{def:coordinate-transform-operator}
\[
\mathcal C_c:
\Omega_{\mathrm{world}}\rightarrow\Omega_{\mathrm{world}}.
\]
\end{definition}

Er ändert die Koordinatenrepräsentation unter einer deklarierten
Transformationsqualifikation. Er bewahrt die Identität des repräsentierten
Objekts und die Provenienz von Quell- und Ziel-CRS. Er kann weder die
World-to-Track-Interpretation noch die metrische Realisierung ersetzen.

\section{Repräsentationen und Beobachtungen erzeugen}

Die Abtastung der überarbeiteten metrischen Trajektorie erzeugt eine
CAD-Kurven- oder Polylinienrepräsentation. Der Repräsentationsoperator
empfängt qualifiziertes Zustandswissen und Repräsentationsparameter; er
erzeugt ein abgeleitetes Artefakt, bewahrt die Objektidentität und hält
Provenienz von Quellzustand und Verfahren fest. Er kann nicht allein dadurch
zur konstruktiven Definition werden, dass er gespeichert, ausgetauscht oder
gerendert wird.

Eine Beobachtungsoperation empfängt dagegen einen physischen Zustand über
einen Messprozess und erzeugt quellengebundene Evidenz mit Zeit, Methode,
Bedingungen und Unsicherheit. Eine Vermessungspolyline ist für den Vergleich
nur durch diese Provenienz nutzbar. Beobachtung kann ihre Ausgabe für sich
genommen weder zur physischen Realität noch zu angenommenem Ingenieurwissen
machen.

\section{Ziel und physischen Zustand verknüpfen}

\begin{definition}[Realisierungsoperator]
\label{def:realization-operator}
\[
\mathcal R:
\mathcal X_{\mathrm{target}}\rightarrow\mathcal X_{\mathrm{real}}.
\]
\end{definition}

Der Operator empfängt einen Zielzustand sowie ein anwendbares Bau-,
Instandhaltungs- oder physisches Antwortmodell und erzeugt die Beschreibung
eines realisierten Zustands. Er bewahrt die Identität des Engineering Objects,
während sich der beschriebene Zustand ändert. Seine Provenienz umfasst Ziel,
Prozessmodell und relevante physische Eingaben. Er kann weder eine Beobachtung
des Ergebnisses erzeugen noch entscheiden, dass der realisierte Zustand
akzeptabel ist.

\section{Die qualifizierte Folge bewerten}

\begin{definition}[Bewertungsoperator]
\label{def:evaluation-operator}
\[
\mathcal E_c:
\mathcal X\times\Omega_{\mathrm{track}}\rightarrow\mathcal Y.
\]
\end{definition}

Für die wiederkehrende Abweichung empfängt die Bewertung den metrisch
realisierten Zielzustand, kompatible Vermessungsevidenz, Vergleichsorte,
Unsicherheit und anwendbare Regeln. Sie erzeugt Residuen oder andere
qualifizierte Folgen. Sie bewahrt Identität und Rollen ihrer Eingaben; ihre
Provenienz hält diese Eingaben, Qualifikationen, Operatorversion und Annahmen
fest. Sie kann weder Evidenz als Wissen annehmen, noch einen Vorschlag
akzeptieren oder eine Ingenieurentscheidung erlassen.

\begin{principle}[Solver-Unabhängigkeit]
\label{princ:operators-solver-independence}
Ingenieurbeobachtungen, Nebenbedingungen, Residuen, Lösungskandidaten und
Entscheidungen sind unabhängig vom numerischen Solver, mit dem sie bewertet
werden.
\end{principle}

\section{Komposition ohne Verantwortungsübertragung}

Für kompatible Typen gilt
\[
\mathcal O_1:\mathcal A\rightarrow\mathcal B,
\qquad
\mathcal O_2:\mathcal B\rightarrow\mathcal C,
\qquad
\mathcal O_2\circ\mathcal O_1:
\mathcal A\rightarrow\mathcal C.
\]

\begin{principle}[Operatortrennung]
\label{princ:operator-separation}
Operatoren mit unterschiedlichen ingenieurtechnischen Verantwortlichkeiten
müssen begrifflich unterscheidbar bleiben, auch wenn sie am selben
Arbeitsablauf beteiligt sind.
\end{principle}

Die Kette der Krümmungsänderung kann somit als Rekonstruktion, metrische
Realisierung, Interpretation einer Repräsentation oder Beobachtung und
Bewertung gelesen werden. Komposition überträgt Werte und Provenienz. Sie
überträgt weder Identität, noch ändert sie stillschweigend den Wahrheitsmodus
oder verleiht der abschließenden Zahl Entscheidungsautorität.
